Before running Ge\+N-\/\+Foam, one has to provide meshes, physical properties, discretization methods (if one does not want to use the default ones) and simulation details. This section provides a quick overview of the input deck of Ge\+N-\/\+Foam. For more details, please refer to the \hyperlink{USERMAN}{User manual}.

\subsection*{Meshing}

Ge\+N-\/\+Foam uses three different meshes for neutronics, thermal-\/hydraulics and thermal-\/mechanics. There is no requirement for the three meshes to occupy the same region of space. Consistent mapping of fields is performed and a reference value is given to a field if no correspondence is found in the mesh where its value is being projected from. It follows that the geometry for neutronics can cover only a small part of the overall reactor geometry. Meshes can be created with every Open\+F\+O\+A\+M-\/compatible tool. Meshes can (should) be divided into zones (cell\+Zones) to allow the use of different physical properties (e.\+g., cross sections) in different reactor regions. Sometimes, when converting a mesh to the Open\+F\+O\+AM (polymesh) format, cell\+Sets (and not cell\+Zones) are created. The topo\+Set\+Dict can be used to convert cell\+Sets to cell\+Zones.

The 3\+D\+\_\+\+Small\+E\+S\+FR tutorial includes an example of mesh generation with Gmsh \cite{https://doi.org/10.1002/nme.2579}. The folder includes the subfolder {\itshape 3d\+E\+S\+F\+R\+Mesh} containing three subfolders for the generation of the meshes for neutronics, thermal-\/hydraulics and thermal-\/mechanics. To create the 3D E\+S\+FR core geometry, one should execute the following command in a terminal\+:

{\ttfamily gmsh esf\+Main.\+geo}

This will create (after a relatively long time) a geometry and open the Gmsh graphical interface. The geometry must then be meshed and the resulting mesh saved and copied in the root of the E\+S\+F3\+D\+\_\+small folder.

By typing in a terminal\+:

{\ttfamily gmsh\+To\+Foam (name of mesh file)}

a {\itshape poly\+Mesh} folder will be created (or updated) in the folder {\itshape constant}. One should then copy this folder in {\itshape constant/neutro\+Region}, {\itshape constant/fluid\+Region} or {\itshape constant/thermo\+Mechanical\+Region}, and repeat the operation for all the three meshes.

Please notice that the 3\+D\+\_\+\+Small\+E\+S\+FR tutorial already contains the correct {\itshape polymesh} folders so that one can avoid the mesh generation step.

N.\+B. {\bfseries A dummy mesh must always be present in all physics (region) directories}, even if not solved for. The E\+M\+P\+TY case is already provided with minimal dummy meshes and consistent fields in the “0” folder. Be careful! In case of parallel calculations all your meshes will have to have a number of cells equal or higher than the number of domains you are decomposing your geometry into. In case you need more cells than what available in the E\+M\+P\+TY case, you can run a refine\+Mesh

\subsection*{Physical properties}

All the data for the Ge\+N-\/\+Foam simulations can be filled in the following input files (dictionaries)\+:
\begin{DoxyItemize}
\item {\itshape constant/thermo\+Mechanical\+Region/thermo\+Mechanical\+Properties} -\/ thermo-\/mechanical properties of structures, subdivided according to the cell\+Zones of the thermo\+Mechanical\+Region mesh. One can find a detailed, commented example in the tutorial \href{https://gitlab.com/foam-for-nuclear/GeN-Foam/-/tree/master/Tutorials/3D_SmallESFR/rootCase/constant/thermoMechanicalRegion/thermoMechanicalProperties}{\tt 3\+D\+\_\+\+Small\+E\+S\+FR}.
\item {\itshape constant/fluid\+Region/g} -\/ gravitational acceleration.
\item {\itshape constant/fluid\+Region/turbulence\+Properties} -\/ standard Open\+F\+O\+AM dictionary to define the turbulence model to be used. One can find a detailed, commented example in the tutorial \href{https://gitlab.com/foam-for-nuclear/GeN-Foam/-/tree/master/Tutorials/3D_SmallESFR/rootCase/constant/fluidRegion/turbulenceProperties}{\tt 3\+D\+\_\+\+Small\+E\+S\+FR}.
\item {\itshape constant/fluid\+Region/thermophysical\+Properties} (for single-\/phase simulations) -\/ standard Open\+F\+O\+AM dictionary to define the thermo-\/physical properties of the coolant. One can find a detailed, commented example in tutorial \href{https://gitlab.com/foam-for-nuclear/GeN-Foam/-/tree/master/Tutorials/3D_SmallESFR/rootCase/constant/fluidRegion/thermophysicalProperties}{\tt 3\+D\+\_\+\+Small\+E\+S\+FR} (single phase)
\item {\itshape constant/fluid\+Region/thermophysical\+Properties.(name of fluid)} (for two-\/phase simulations) -\/ standard Open\+F\+O\+AM dictionaries to define the thermo-\/physical properties of various phases. The name of fluid is defined in {\itshape constant/fluid\+Region/phase\+Properties}. One can find a detailed, commented example in the tutorial \href{https://gitlab.com/foam-for-nuclear/GeN-Foam/-/tree/master/Tutorials/1D_boiling/constant/fluidRegion/thermophysicalProperties.liquid}{\tt 1\+D\+\_\+boiling (liquid)}, \href{https://gitlab.com/foam-for-nuclear/GeN-Foam/-/tree/master/Tutorials/1D_boiling/constant/fluidRegion/thermophysicalProperties.vapour}{\tt (vapour)}.
\item {\itshape constant/fluid\+Region/phase\+Properties} -\/ large dictionary that can be used to\+: determined whether the simulation is single-\/phase or two-\/phase; set various properties of the phases (beside the thermo-\/physical properties defined in {\itshape constant/fluid\+Region/thermophysical\+Properties}); set the properties of the sub-\/scale structures (fuel pins, heat exchangers, etc) in the porous zones, including the possibility to assign a {\itshape power\+Model} for power production (e.\+g., nuclear fuel, or constant power) and the {\itshape passive\+Properties} of another sub-\/structure that interacts thermally with the fluid (for instance the wrappers in sodium fast reactors). The name of the porous zones must coincide with that of the cell\+Zones of the fluid\+Region mesh. Anisotropic pressure drops can be set by using the keywords {\itshape transverse\+Drag\+Model} (Blasius, Gunter\+Shaw, same) and {\itshape principal\+Axis}(localX, localY, localZ) in the sub-\/dictionary {\itshape drag\+Models.(name\+Of\+Phase).structure.(name\+Of\+Cell\+Zones)}. {\itshape principal\+Axis} sets the axis on which the nominal drag\+Model is used. {\itshape transverse\+Drag\+Model} sets the model to be used on the two directions that are perpendicular to {\itshape principal\+Axis}. If {\itshape same} is chosen as {\itshape transverse\+Drag\+Model}, the code will use the nominal model in all directions, but with the possibility of an anisotropic hydraulic diameter. The anisotropy of the hydraulic diameter can be set using the keyword {\itshape local\+Dh\+Anisotrpy} and assign to it a vector of 3 scaling factors (one for each local directions). One can find detailed, commented examples in the tutorials \href{https://gitlab.com/foam-for-nuclear/GeN-Foam/-/tree/master/Tutorials/3D_SmallESFR/rootCase/constant/fluidRegion/phaseProperties}{\tt 3\+D\+\_\+\+Small\+E\+S\+FR} (single phase) and \href{https://gitlab.com/foam-for-nuclear/GeN-Foam/-/tree/master/Tutorials/1D_boiling/constant/fluidRegion/phaseProperties}{\tt 1\+D\+\_\+boiling} (two phases).
\item {\itshape constant/neutro\+Region/neutronics\+Properties} -\/ dictionary to control how neutronics is solved (point kinetics, diffusion, S\+P3 or SN), and if it\textquotesingle{}s an eigenvalue calculation or a transient. One can find detailed, commented examples in most tutorials. See for instance \href{https://gitlab.com/foam-for-nuclear/GeN-Foam/-/tree/master/Tutorials/3D_SmallESFR/rootCase/constant/neutroRegion/neutronicsProperties}{\tt 3\+D\+\_\+\+Small\+E\+S\+FR} (single phase).
\item {\itshape constant/neutro\+Region/reactor\+State} – contains the target power (p\+Target) for eigenvalue calculations, the keff that results from the eigenvalue calculations and the external reactivity (i.\+e., the extra reactivity one can add for instance to simulate a reactivity step). N.\+B.\+: keff has no effect on point\+Kinetics. You can find detailed, commented examples in most tutorials. N.\+B.\+2\+: In point kinetics, p\+Target is the initial value used by the point kinetics solver to plot results, but the solver actually scale the power\+Density and flux fields provided by the user. It is up to the user to make sure that p\+Target is consistent with the power\+Density and flux fields. A commented reactor\+State can be found in \href{https://gitlab.com/foam-for-nuclear/GeN-Foam/-/tree/master/Tutorials/3D_SmallESFR/rootCase/constant/neutroRegion/reactorState}{\tt 3\+D\+\_\+\+Small\+E\+S\+FR} (single phase). Please note that eigenvalue calculations will update the keff value in this dictionary. In parallel calculations, the updated value can be found in {\itshape processor0/constant/neutro\+Region/reactor\+State}.
\item {\itshape constant/neutro\+Region/nuclear\+Data} -\/ contains all basic nuclear properties for the reference reactor state. The other {\itshape nuclear\+Data...} files in {\itshape constant/neutronics/} should include the cross-\/sections for perturbed reactor states. In addition, these files include information about the perturbed and reference ({\itshape nuclear\+Data}) reactor state. For instance, {\itshape nuclear\+Data\+Fuel\+Temp} must include {\itshape Tfuel\+Ref} and {\itshape Tfuel\+Perturbed}, which represent the temperatures at which the reference ({\itshape nuclear\+Data}) and perturbed ({\itshape nuclear\+Data\+Fuel\+Temp}) cross sections have been calculated, respectively. Linear interpolation is performed by Ge\+N-\/\+Foam between reference and perturbed reactor states, except for fuel temperature, for which a logarithmic or square root interpolation is provided (depending on the spectrum, which in turns is defined by the keyword {\itshape fast\+Neutrons}). If no data are provided, the reference cross sections are used. Nuclear data can be generated using any nuclear code. The \href{https://gitlab.com/foam-for-nuclear/GeN-Foam/-/tree/master/Tools/serpentToFoam/serpent2.1.23}{\tt serpent\+To\+Foam} routines provided with Ge\+N-\/\+Foam (in the {\itshape Tools} folder) is an Octave script that automatically converts Serpent output files into the nuclear data files employed by Ge\+N-\/\+Foam. The entry {\itshape disc\+Factor} is used only if discontinuity factors have to be used. The term {\itshape integral\+Flux}, is used only if the automatic adjustment of discontinuity factors is performed \cite{FIORINA2016212}. Nonetheless, these entries should always be present. One can find detailed, commented examples of nuclear\+Data in the tutorials \href{https://gitlab.com/foam-for-nuclear/GeN-Foam/-/tree/master/Tutorials/3D_SmallESFR/rootCase/constant/neutroRegion/nuclearData}{\tt 3\+D\+\_\+\+Small\+E\+S\+FR} (for diffusion or S\+P3), \href{https://gitlab.com/foam-for-nuclear/GeN-Foam/-/tree/master/Tutorials/Godiva_SN/constant/neutroRegion/nuclearData}{\tt Godiva\+\_\+\+SN} (for discrete ordinates) and \href{https://gitlab.com/foam-for-nuclear/GeN-Foam/-/tree/master/Tutorials/2D_onePhaseAndPointKineticsCoupling/rootCase/constant/neutroRegion/nuclearData}{\tt 2\+D\+\_\+one\+Phase\+And\+Point\+Kinetics\+Coupling} (for point kinetics). One can find examples of the {\itshape nuclear\+Data...} files in the tutorial \href{https://gitlab.com/foam-for-nuclear/GeN-Foam/-/tree/master/Tutorials/3D_SmallESFR/rootCase/constant/neutroRegion}{\tt 3\+D\+\_\+\+Small\+E\+S\+FR}
\item {\itshape constant/neutro\+Region/quadrature\+Set} -\/ contains the quadrature set for discrete ordinate calculations. One can find examples of three different quadrature set in the tutorial \href{https://gitlab.com/foam-for-nuclear/GeN-Foam/-/tree/master/Tutorials/Godiva_SN/constant/neutroRegion/}{\tt Godiva\+\_\+\+SN}.
\item {\itshape constant/neutro\+Region/\+C\+R\+Move} -\/ contains input data for control rods movement. Control rods can be moved from the initial position to a new one by selecting initial and final time of the insertion/extraction and the speed of insertion/extraction (positive speed for insertion). One can find a commented example in the tutorial \href{https://gitlab.com/foam-for-nuclear/GeN-Foam/-/tree/master/Tutorials/3D_SmallESFR/rootCase/constant/neutroRegion/CRmove}{\tt 3\+D\+\_\+\+Small\+E\+S\+FR}, though this option is not actually used in the tutorial.
\end{DoxyItemize}

\subsection*{Initial values and boundary conditions}

As in all standard Open\+F\+O\+AM solvers, initial values (IC) and boundary conditions (BC) should be provided in the “0” folder, or in the time folder corresponding to the {\itshape start\+Time} of the simulation, if different than 0.

\subsection*{Discretization and solution}

Details for discretization and solution of single-\/physics equations are handled in a standard Open\+F\+O\+AM way, i.\+e., through the {\itshape fv\+Solution} and {\itshape fv\+Schemes} dictionaries in {\itshape constant/neutro\+Region}, {\itshape constant/fluid\+Region} or {\itshape constant/ thermal\+Mechanical\+Region}. Coupling and simulation details are determined through the {\itshape fv\+Solution} and {\itshape control\+Dict} in the {\itshape system} folder. The {\itshape control\+Dict} is significantly extended compared to a standard Open\+F\+O\+AM control\+Dict in order to control what solvers are activated and flags that affect the behaviour of Ge\+N-\/\+Foam as a whole. 